\chapter*{Lesepfad und Beweiszuordnung}
\addcontentsline{toc}{chapter}{Lesepfad und Beweiszuordnung}
Die \href{_B19-cuts-reading.pdf\#cuts.reading}{Lesefassung}
führt mit kurzen Beweisen von rationalen Hauptschnitten bis zum Supremum.
Das \href{_B19-cuts-exercises.pdf\#cuts.exercises}{Übungsheft}
enthält Aufgaben und Lösungen. Dieser Band ergänzt die vollständigen
Tabellen des mathematischen Kerns. Sämtliche Aussagen behalten ihre Nummern
und ihre Referenzziele in \href{../_B19.pdf\#Doc-Start}{Band~19}.
Die hier gedruckten Tabellen und die Tabellen im Fachband werden aus
\emph{denselben Quelldateien} erzeugt.

\begin{center}
\small\renewcommand{\arraystretch}{1.2}
\begin{tabularx}{\linewidth}{@{}p{.38\linewidth}X@{}}
\toprule
Schritt der Lesefassung & Kanonische Begründung\\
\midrule
Schnittbedingungen & \FormulaRefAuto{DedekindCutPredicate}[def] und
\FormulaRefAuto{DedekindCutCharacterization}[thm]\\
Rationale Hauptschnitte & \FormulaRefAuto{RealPrincipalCutIsCut}[thm]\\
Vergleich durch Inklusion & \FormulaRefAuto{RealCutInclusionComparable}[thm]\\
Ordnungstreue und Injektivität & \FormulaRefAuto{RationalToRealEmbeddingOrder}[thm] und
\FormulaRefAuto{RationalToRealEmbeddingInjective}[thm]\\
Die Vereinigung ist ein Schnitt & \FormulaRefAuto{RealBoundedCutUnionIsCut}[thm]\\
Vollständigkeit und Supremum & \FormulaRefAuto{RealLeastUpperBoundProperty}[thm] und
\FormulaRefAuto{RealCutSupremumUnion}[thm]\\
\bottomrule
\end{tabularx}
\end{center}

Die Beispiele wenden diese Aussagen auf konkrete Mengen an. Die kurzen
Lesebeweise wählen bei Hauptschnitten auch den Randwert als Gegenzeugen;
die Tabellen führen den Ordnungsvergleich über die rationale Dichtheit aus.
Vorausgesetzt werden rationale Arithmetik und Ordnung (Band~18),
Mengenoperationen (Band~03) und der Supremumsbegriff (Band~13).

\paragraph{Korrigierte Voraussetzungen.}
\FormulaRefAuto{RealCutUnionUpperBound}[thm] setzt nun ausdrücklich
$S\in\mathbb R$ voraus, \FormulaRefAuto{RealCutUnionLeastUpperBound}[thm]
entsprechend $\bigcup\mathcal A\in\mathbb R$. Beide Anwendungen erhalten
diesen Nachweis aus \FormulaRefAuto{RealBoundedCutUnionIsCut}[thm];
ihre Abhängigkeitslisten führen die Schrankenannahme mit.

\clearpage
\setcounter{chapter}{1}
% Kanonische Quelle: Band 19 und Beweistabellen verwenden denselben Text.
\chapter{Dedekindsche Schnitte}

\section{Schnittprädikat und Trägermenge}

\FormulaDefDeltaK[Dedekindscher Schnitt der rationalen Zahlen]{%
  \RealCut{A}%
}{DedekindCutPredicate}{%
  \DeltaRow{Mengen}{A\dsep p\dsep q\dsep r}
  \DeltaRow{Relationsoperatoren}{<}
  \DeltaPrem{Geordneter rationaler Körper}{\mathbb Q}
  \DeltaRow{\textbf{Axiome}}{}
  \DeltaPrem{Rationale Teilmenge}{A\subseteq\mathbb Q}
  \DeltaPrem{Nichtleerheit}{A\neq\varnothing}
  \DeltaPrem{Echtheit}{A\neq\mathbb Q}
  \DeltaRow{Nach unten abgeschlossen}{%
    p\in A\dsep q\in\mathbb Q\dsep q<p\vdash q\in A}
  \DeltaRow{Kein größtes Element}{p\in A\vdash\exists r\in A\,(p<r)}
  \DeltaRow{\textbf{Neues Symbol}}{}
  \DeltaPrem{Schnittprädikat}{\RealCut{A}}
}

\FormulaAxiomDeltaK[Rationale Trägermenge]{%
  \RealCut{A}\vdash A\subseteq\mathbb Q%
}{DedekindCutSubsetAxiom}{%
  \DeltaRow{Mengen}{A\dsep p\dsep q\dsep r}
}

\FormulaAxiomDeltaK[Nichtleerheit]{%
  \RealCut{A}\vdash A\neq\varnothing%
}{DedekindCutNonemptyAxiom}{%
  \DeltaRow{Mengen}{A\dsep p\dsep q\dsep r}
}

\FormulaAxiomDeltaK[Echter Schnitt]{%
  \RealCut{A}\vdash A\neq\mathbb Q%
}{DedekindCutProperAxiom}{%
  \DeltaRow{Mengen}{A\dsep p\dsep q\dsep r}
}

\FormulaAxiomDeltaK[Abwärtsabgeschlossenheit]{%
  \RealCut{A}\vdash \forall p\in A\,\forall q\in\mathbb Q\,(q<p\rightarrow q\in A)%
}{DedekindCutDownwardAxiom}{%
  \DeltaRow{Mengen}{A\dsep p\dsep q\dsep r}
}

\FormulaAxiomDeltaK[Kein größtes Element]{%
  \RealCut{A}\vdash \forall p\in A\,\exists r\in A\,(p<r)%
}{DedekindCutNoMaximumAxiom}{%
  \DeltaRow{Mengen}{A\dsep p\dsep q\dsep r}
}

\FormulaThmDeltaK[Schnittkriterium]{%
  \begin{aligned}[t]
  \RealCut{A}\eqvdash{}&
    A\subseteq\mathbb Q\land A\neq\varnothing\land A\neq\mathbb Q
    \land{}\\[-2pt]
  &\forall p\in A\,\forall q\in\mathbb Q\,
      \bigl(q<p\rightarrow q\in A\bigr)
    \land{}\\[-2pt]
  &\forall p\in A\,\exists r\in A\,(p<r).
  \end{aligned}%
}{DedekindCutCharacterization}{%
  \DeltaRow{Mengen}{A\dsep p\dsep q\dsep r}
  \DeltaRow{Relationsoperatoren}{<}
}
\begin{tabproofsplitwide}
  \proofpartwide{Hinrichtung}
    \proofstepwidestar[1]{\RealCut{A}}{\rA}
    \proofstepwidestar[1]{A\subseteq\mathbb Q}{%
      \FormulaRefAuto{DedekindCutSubsetAxiom}[ax]{1}}
    \proofstepwidestar[1]{A\neq\varnothing}{%
      \FormulaRefAuto{DedekindCutNonemptyAxiom}[ax]{1}}
    \proofstepwidestar[1]{A\neq\mathbb Q}{%
      \FormulaRefAuto{DedekindCutProperAxiom}[ax]{1}}
    \proofstepwidestar[1]{\forall p\in A\,\forall q\in\mathbb Q\,(q<p\rightarrow q\in A)}{%
      \FormulaRefAuto{DedekindCutDownwardAxiom}[ax]{1}}
    \proofstepwidestar[1]{\forall p\in A\,\exists r\in A\,(p<r)}{%
      \FormulaRefAuto{DedekindCutNoMaximumAxiom}[ax]{1}}
    \proofstepwidestar[1]{%
      \begin{aligned}[t]
      &A\subseteq\mathbb Q\land A\neq\varnothing\land A\neq\mathbb Q\land{}\\[-2pt]
      &\forall p\in A\,\forall q\in\mathbb Q\,(q<p\rightarrow q\in A)\land{}\\[-2pt]
      &\forall p\in A\,\exists r\in A\,(p<r)
    \end{aligned}}{\rAIStar{2,3,4,5,6}}
  \closeproofpartwide

  \proofpartwide{Rückrichtung}
    \proofstepwidestar[1]{%
      \begin{aligned}[t]
      &A\subseteq\mathbb Q\land A\neq\varnothing\land A\neq\mathbb Q\land{}\\[-2pt]
      &\forall p\in A\,\forall q\in\mathbb Q\,(q<p\rightarrow q\in A)\land{}\\[-2pt]
      &\forall p\in A\,\exists r\in A\,(p<r)
    \end{aligned}}{\rA}
    \proofstepwidestar[1]{A\subseteq\mathbb Q}{\rAEa{1}}
    \proofstepwidestar[1]{A\neq\varnothing}{\rAEa{\rAEb{1}}}
    \proofstepwidestar[1]{A\neq\mathbb Q}{\rAEa{\rAEb{\rAEb{1}}}}
    \proofstepwidestar[1]{\forall p\in A\,\forall q\in\mathbb Q\,(q<p\rightarrow q\in A)}{\rAEa{\rAEb{\rAEb{\rAEb{1}}}}}
    \proofstepwidestar[1]{\forall p\in A\,\exists r\in A\,(p<r)}{\rAEb{\rAEb{\rAEb{\rAEb{1}}}}}
    \proofstepwidestar[1]{\RealCut{A}}{%
      \FormulaRefAuto{DedekindCutPredicate}[def]{2,3,4,5,6}}
  \closeproofpartwide
\end{tabproofsplitwide}

Neben der rationalen Typisierung sind die vier Schnittbedingungen wesentlich. Die
Abwärtsabgeschlossenheit macht einen Schnitt zu einem rationalen
Anfangsabschnitt; die letzte Bedingung entfernt die ansonsten mögliche
Mehrdeutigkeit am Rand eines rationalen Anfangsabschnitts.

\FormulaDefDeltaK[Die Menge der reellen Zahlen]{%
  \mathbb R\coloneqq
  \bigl\{A\in\powerset(\mathbb Q)\bigm|\RealCut{A}\bigr\}%
}{RealCutCarrier}{%
  \DeltaRow{Mengen}{A\dsep\mathbb R}%
  \DeltaRow{Prädikate}{\RealCut{\,\cdot\,}}%
  \DeltaRow{\textbf{Neues Symbol}}{}%
  \DeltaRow{Reelle Zahlen}{\mathbb R}%
}

\FormulaThmDeltaK[Elementcharakterisierung der reellen Zahlen]{%
  A\in\mathbb R\eqvdash\RealCut{A}%
}{RealCutMembership}{%
  \DeltaRow{Mengen}{A}%
}
\begin{tabproofsplitwide}
  \proofpartwideindR[Hinrichtung]{%
    A\in\mathbb R\rightarrow\RealCut A%
  }{A\in\mathbb R\rightarrow\RealCut A}[RealCutMembershipForward]
    \proofstepwidestar[1]{A\in\mathbb R}{\rA}
    \proofstepwidestar[]{%
      \mathbb R=\{X\in\powerset(\mathbb Q)\mid\RealCut X\}}{%
      \FormulaRefAuto{RealCutCarrier}[def]}
    \proofstepwidestar[1]{%
      A\in\powerset(\mathbb Q)\land\RealCut A}{%
      \rIE{1,2}}
    \proofstepwidestar[1]{\RealCut A}{\rAEb{3}}
  \proofstepwidestar[]{A\in\mathbb R\rightarrow\RealCut A}{\rRI{1,4}}
  \closeproofpartwideindR

  \proofpartwideindR[Rückrichtung]{%
    \RealCut A\rightarrow A\in\mathbb R%
  }{\RealCut A\rightarrow A\in\mathbb R}[RealCutMembershipBackward]
    \proofstepwidestar[1]{\RealCut A}{\rA}
    \proofstepwidestar[1]{A\subseteq\mathbb Q}{%
      \FormulaRefAuto{DedekindCutCharacterization}[thm]{1}}
    \proofstepwidestar[1]{A\in\powerset(\mathbb Q)}{%
      \FormulaRefAuto{x\in\powerset(A)\eqvdash x\subseteq A}{2}}
    \proofstepwidestar[1]{%
      A\in\powerset(\mathbb Q)\land\RealCut A}{\rAI{3,1}}
    \proofstepwidestar[1]{A\in\mathbb R}{%
      \FormulaRefAuto{RealCutCarrier}[def]{4}}
  \proofstepwidestar[]{\RealCut A\rightarrow A\in\mathbb R}{\rRI{1,5}}
  \closeproofpartwideindR

  \proofpartwide{Zusammenfassung}
    \proofstepwidestar[]{A\in\mathbb R\rightarrow\RealCut A}{%
      \FormulaRefAuto{RealCutMembershipForward}[thm]}
    \proofstepwidestar[]{\RealCut A\rightarrow A\in\mathbb R}{%
      \FormulaRefAuto{RealCutMembershipBackward}[thm]}
    \proofstepwidestar[]{A\in\mathbb R\leftrightarrow\RealCut A}{%
      \rLRI{1,2}}
  \closeproofpartwide
\end{tabproofsplitwide}

Diese Charakterisierung erlaubt es im Folgenden, eine reelle Zahl und den sie
darstellenden Schnitt ohne zusätzliche Klassenklammern miteinander zu
identifizieren.

\section{Hauptschnitte und rationale Einbettung}

\FormulaDefDeltaK[Hauptschnitt einer rationalen Zahl]{%
  q\in\mathbb Q\vdash
  \RealHat q\coloneqq\{p\in\mathbb Q\mid p<q\}%
}{RealPrincipalCut}{%
  \DeltaRow{Mengen}{p\dsep q}%
  \DeltaRow{Relationsoperatoren}{<}%
  \DeltaRow{\textbf{Neues Symbol}}{}%
  \DeltaRow{Hauptschnitt}{\RealHat q}%
}

\FormulaThmDeltaK[Hauptschnitte sind Dedekindsche Schnitte]{%
  q\in\mathbb Q\vdash\RealCut{\RealHat q}%
}{RealPrincipalCutIsCut}{%
  \DeltaRow{Mengen}{p\dsep q\dsep r}%
}
\begin{tabproofsplitwide}
\proofpartwide{Nichtleerheit}
\proofstepwidestarbreakkeep[1]{q\in\mathbb Q}{\rA}
    \proofstepwidestarbreakkeep[]{\IntOne\in\PositiveIntegers}{%
      \FormulaRefAuto{IntegerOnePositive}[thm]}
    \proofstepwidestarbreakkeep[]{\IntZero<\IntOne}{%
      \FormulaRefAuto{PositiveIntegersDef}[def]{2}}
    \proofstepwidestarbreakkeep[]{%
      \IntZero\leq\IntOne\land\IntZero\neq\IntOne}{%
      \FormulaRefAuto{IntegerStrictOrderDef}[def]{3}}
    \proofstepwidestarbreakkeep[]{\IntZero\in\mathbb Z}{%
      \FormulaRefAuto{IntZeroCarrier}[thm]}
    \proofstepwidestarbreakkeep[]{\IntOne\in\mathbb Z}{%
      \FormulaRefAuto{IntOneCarrier}[thm]}
    \proofstepwidestarbreakkeep[]{%
      \RatEmbed(\IntZero)\leq\RatEmbed(\IntOne)}{%
      \FormulaRefAuto{IntegerRationalEmbeddingOrder}[thm]{4,5,6}}
    \proofstepwidestarbreakkeep[]{%
      \RatEmbed(\IntZero)\neq\RatEmbed(\IntOne)}{%
      \FormulaRefAuto{IntegerRationalEmbeddingInjective}[thm]{4,5,6}}
    \proofstepwidestar[]{\RatOne=\RatEmbed(\IntOne)}{%
      \FormulaRefAuto{RationalOneDef}[def]}
    \proofstepwidestar[]{\RatZero=\RatEmbed(\IntZero)}{%
      \FormulaRefAuto{RationalZeroDef}[def]}
    \proofstepwidestarbreakkeep[]{%
      \RatZero=\RatEmbed(\IntZero)\land
      \RatOne=\RatEmbed(\IntOne)}{%
      \rAI{10,
        9}}
    \proofstepwidestarbreakkeep[]{\RatZero<\RatOne}{%
      \FormulaRefAuto{RationalStrictOrderDef}[def]{7,8,11}}
    \proofstepwide[1,12]{q-\RatOne}{<}{q}{%
      \FormulaRefAuto{RationalOrderTranslation}[thm]{1,12}}
    \proofstepwidestarbreakkeep[1,12]{%
      \exists p\in\mathbb Q\,(q-\RatOne<p\land p<q)}{%
      \FormulaRefAuto{RationalOrderDense}[thm]{13}}
    \proofstepwidestarbreakkeep[1,12]{p\in\mathbb Q\land p<q}{\rA}
    \proofstepwidestarbreakkeep[1,12]{p\in\RealHat q}{%
      \FormulaRefAuto{RealPrincipalCut}[def]{15}}
    \proofstepwidestarbreakkeep[1,12]{\RealHat q\neq\varnothing}{%
      \FormulaRefAuto{x\in A\vdash A\neq\varnothing}{16}}

\closeproofpartwide

\proofpartwide{Echtheit}
  \setcounter{proofstepnr}{17}% Fortlaufende Schritte dieses Beweises.
\proofstepwidestar[1]{\TotOrd{\mathbb Q,\leq}}{%
      \FormulaRefAuto{RationalOrderTotal}[thm]{1}}
\proofstepwidestarbreakkeep[1]{q\notin\RealHat q}{%
      \FormulaRefAuto{RealPrincipalCut}[def]{%
        18}}
    \proofstepwidestarbreakkeep[1]{\RealHat q\neq\mathbb Q}{%
      \FormulaRefAuto{x\in B, x\not\in A\vdash A\neq B}{1,19}}

\closeproofpartwide

\proofpartwideindR[Nichtleerheit und Echtheit]{%
    q\in\mathbb Q\vdash
    \RealHat q\neq\varnothing\land\RealHat q\neq\mathbb Q
  }{%
    q\in\mathbb Q\vdash
    \RealHat q\neq\varnothing\land\RealHat q\neq\mathbb Q
  }[RealPrincipalCutNonemptyProper]
  \setcounter{proofstepnr}{20}% Fortlaufende Schritte dieses Beweises.
\proofstepwidestarbreakkeep[1]{%
      \RealHat q\neq\varnothing\land\RealHat q\neq\mathbb Q}{%
      \rAI{17,20}}
    \proofstepwidestarbreakkeep[1]{%
      \RealHat q\neq\varnothing\land\RealHat q\neq\mathbb Q}{%
      \rEE{14,15,21}}

\closeproofpartwideindR


  \proofpartwideindR[Abwärtsabschluss]{%
    \begin{aligned}[t]
      &q\in\mathbb Q\dsep p\in\RealHat q\dsep s\in\mathbb Q\dsep s<p\\[-2pt]
      &\quad\vdash s\in\RealHat q
    \end{aligned}
  }{%
    q\in\mathbb Q\dsep p\in\RealHat q\dsep s\in\mathbb Q\dsep s<p
    \vdash s\in\RealHat q
  }[RealPrincipalCutDownwardClosed]
    \proofstepwidestar[1]{q\in\mathbb Q}{\rA}
    \proofstepwidestar[2]{p\in\RealHat q}{\rA}
    \proofstepwidestar[3]{s\in\mathbb Q}{\rA}
    \proofstepwidestar[4]{s<p}{\rA}
    \proofstepwide[1,2]{p}{<}{q}{%
      \FormulaRefAuto{RealPrincipalCut}[def]{1,2}}
    \proofstepwide[1,2,4]{s}{<}{q}{%
      \FormulaRefAuto{RationalOrderTotal}[thm]{4,5}}
    \proofstepwidestar[1,2,3,4]{s\in\RealHat q}{%
      \FormulaRefAuto{RealPrincipalCut}[def]{3,6}}
  \closeproofpartwideindR

  \proofpartwideindR[Kein größtes Element]{%
    q\in\mathbb Q\dsep p\in\RealHat q
    \vdash\exists r\in\RealHat q\,(p<r)
  }{%
    q\in\mathbb Q\dsep p\in\RealHat q
    \vdash\exists r\in\RealHat q\,(p<r)
  }[RealPrincipalCutNoMaximum]
    \proofstepwidestar[1]{q\in\mathbb Q}{\rA}
    \proofstepwidestar[2]{p\in\RealHat q}{\rA}
    \proofstepwide[1,2]{p}{<}{q}{%
      \FormulaRefAuto{RealPrincipalCut}[def]{1,2}}
    \proofstepwidestar[1,2]{%
      \exists r\in\mathbb Q\,(p<r\land r<q)}{%
      \FormulaRefAuto{RationalOrderDense}[thm]{3}}
    \proofstepwidestar[1,2]{r\in\mathbb Q\land p<r\land r<q}{\rA}
    \proofstepwidestar[1,2]{r\in\RealHat q}{%
      \FormulaRefAuto{RealPrincipalCut}[def]{5}}
    \proofstepwidestar[1,2]{\exists r\in\RealHat q\,(p<r)}{%
      \rEI{5,6}}
    \proofstepwidestar[1,2]{\exists r\in\RealHat q\,(p<r)}{%
      \rEE{4,5,7}}
  \closeproofpartwideindR

  \proofpartwide{Schluss}
    \proofstepwidestar[1]{q\in\mathbb Q}{\rA}
    \proofstepwidestar[1]{\RealHat q\subseteq\mathbb Q}{%
      \FormulaRefAuto{RealPrincipalCut}[def]{1}}
    \proofstepwidestar[1]{%
      \RealHat q\neq\varnothing\land\RealHat q\neq\mathbb Q}{%
      \FormulaRefAuto{RealPrincipalCutNonemptyProper}[thm]{1}}
    \proofstepwidestar[1]{%
      \forall p\in\RealHat q\,\forall s\in\mathbb Q\,
      (s<p\rightarrow s\in\RealHat q)}{%
      \FormulaRefAuto{RealPrincipalCutDownwardClosed}[thm]{1}}
    \proofstepwidestar[1]{%
      \forall p\in\RealHat q\,\exists r\in\RealHat q\,(p<r)}{%
      \FormulaRefAuto{RealPrincipalCutNoMaximum}[thm]{1}}
    \proofstepwidestar[1]{\RealCut{\RealHat q}}{%
      \FormulaRefAuto{DedekindCutCharacterization}[thm]{2,3,4,5}}
  \closeproofpartwide
\end{tabproofsplitwide}

\FormulaDefDeltaK[Kanonische Einbettung der rationalen Zahlen]{%
  \begin{aligned}[t]
    &\RealEmb\colon\mathbb Q\to\mathbb R,\\[-2pt]
    &\forall q\in\mathbb Q\,
      \bigl(\RealEmb(q)\coloneqq\RealHat q\bigr)
  \end{aligned}%
}{RationalToRealEmbedding}{%
  \DeltaRow{Mengen}{p\dsep q}%
  \DeltaRow{Funktionen}{\RealEmb}%
  \DeltaRow{\textbf{Neues Symbol}}{}%
  \DeltaRow{Rationale Einbettung}{\RealEmb}%
}
\begin{tabproof}
  \proofstep{1}{q\in\mathbb Q}{\rA}
  \proofstep{1}{\RealCut{\RealHat q}}{%
    \FormulaRefAuto{RealPrincipalCutIsCut}[thm]{1}}
  \proofstep{1}{\RealHat q\in\mathbb R}{%
    \FormulaRefAuto{A\in\mathbb R\eqvdash\RealCut{A}}[thm]{2}}
  \proofstep{}{\forall q\in\mathbb Q\,\bigl(\RealHat q\in\mathbb R\bigr)}{%
    \rUI{\rRI{1,3}}}
\end{tabproof}

\FormulaThmDeltaK[Die rationale Einbettung ist eine Funktion]{%
  \RealEmb\colon\mathbb Q\to\mathbb R%
}{RationalToRealEmbeddingFunction}{%
  \DeltaRow{Funktionen}{\RealEmb}%
}
\begin{tabproof}
  \proofstep{}{\RealEmb\colon\mathbb Q\to\mathbb R}{%
    \FormulaRefAuto{RationalToRealEmbedding}[def]}
\end{tabproof}


% Kanonische Quelle: Band 19 und Beweistabellen verwenden denselben Text.
\chapter{Ordnung und Vollständigkeit}

\section{Inklusionsordnung}

\FormulaDefDeltaK[Ordnung der reellen Zahlen]{%
  A,B\in\mathbb R\vdash
  A\RealLe B\coloneqq A\subseteq B%
}{RealCutOrder}{%
  \DeltaRow{Mengen}{A\dsep B}%
  \DeltaRow{Relationsoperatoren}{\RealLe}%
  \DeltaRow{\textbf{Neues Symbol}}{}%
  \DeltaRow{Reelle Ordnung}{\RealLe}%
}

\FormulaDefDeltaK[Strikte Ordnung der reellen Zahlen]{%
  A,B\in\mathbb R\vdash
  A\RealLt B\coloneqq A\subsetneq B%
}{RealCutStrictOrder}{%
  \DeltaRow{Mengen}{A\dsep B}%
  \DeltaRow{Relationsoperatoren}{\RealLe\dsep\RealLt}%
  \DeltaRow{\textbf{Neues Symbol}}{}%
  \DeltaRow{Strikte reelle Ordnung}{\RealLt}%
}

\FormulaThmDeltaK[Charakterisierung der reellen Vergleichsrelationen]{%
  A,B\in\mathbb R\vdash
  \bigl(A\RealLe B\leftrightarrow A\subseteq B\bigr)
  \land
  \bigl(A\RealLt B\leftrightarrow A\subsetneq B\bigr)%
}{RealCutOrderCharacterization}{%
  \DeltaRow{Mengen}{A\dsep B}%
  \DeltaRow{Relationsoperatoren}{\RealLe\dsep\RealLt}%
}
\begin{tabproofwide}
  \proofstepwidestar[1]{A,B\in\mathbb R}{\rA}
  \proofstepwide[1]{A\RealLe B}{\leftrightarrow}{A\subseteq B}{%
    \FormulaRefAuto{RealCutOrder}[def]{1}}
  \proofstepwide[1]{A\RealLt B}{\leftrightarrow}{A\subsetneq B}{%
    \FormulaRefAuto{RealCutStrictOrder}[def]{1}}
  \proofstepwidestar[1]{%
    (A\RealLe B\leftrightarrow A\subseteq B)\land
    (A\RealLt B\leftrightarrow A\subsetneq B)}{\rAI{2,3}}
\end{tabproofwide}

\FormulaThmDeltaK[Vergleichbarkeit zweier Dedekindscher Schnitte]{%
  A,B\in\mathbb R\vdash A\subseteq B\lor B\subseteq A%
}{RealCutInclusionComparable}{%
  \DeltaRow{Mengen}{A\dsep B\dsep a\dsep b}%
}
\begin{tabproofsplitwide}
  \proofpartwideindR[Ausschluss zweier Gegenzeugen]{%
    \begin{aligned}[t]
      &A,B\in\mathbb R\dsep a\in A\setminus B\dsep b\in B\setminus A\\[-2pt]
      &\quad\vdash\bot
    \end{aligned}
  }{%
    A,B\in\mathbb R\dsep a\in A\setminus B\dsep b\in B\setminus A
    \vdash\bot
  }[RealCutOppositeWitnessesImpossible]
    \proofstepwidestar[1]{A,B\in\mathbb R}{\rA}
    \proofstepwidestar[2]{a\in A\setminus B}{\rA}
    \proofstepwidestar[3]{b\in B\setminus A}{\rA}
    \proofstepwidestarbreakkeep[1]{\RealCut A\land\RealCut B}{%
      \FormulaRefAuto{RealCutMembership}[thm]{1}}
    \proofstepwidestar[2]{a\in A\land a\notin B}{%
      \FormulaRefAuto{x\in A\setminus B\eqvdash x\in A\land x\notin B}{2}}
    \proofstepwidestar[3]{b\in B\land b\notin A}{%
      \FormulaRefAuto{x\in A\setminus B\eqvdash x\in A\land x\notin B}{3}}
    \proofstepwidestar[1,2,3]{a<b\lor a=b\lor b<a}{%
      \FormulaRefAuto{RationalOrderTotal}[thm]{1,5,6}}
    \proofstepwidestar[8]{a<b}{\rA}
    \proofstepwidestar[1,2,3,8]{a\in B}{%
      \FormulaRefAuto{DedekindCutCharacterization}[thm]{4,6,8}}
    \proofstepwidestar[1,2,3,8]{\bot}{\rBI{5,9}}
    \proofstepwidestar[11]{a=b}{\rA}
    \proofstepwidestar[1,2,3,11]{a\in B}{\rIE{11,6}}
    \proofstepwidestar[1,2,3,11]{\bot}{\rBI{5,12}}
    \proofstepwidestar[14]{b<a}{\rA}
    \proofstepwidestar[1,2,3,14]{b\in A}{%
      \FormulaRefAuto{DedekindCutCharacterization}[thm]{4,5,14}}
    \proofstepwidestar[1,2,3,14]{\bot}{\rBI{6,15}}
    \proofstepwidestar[1,2,3]{\bot}{\rOE{7,8,10,11,13,14,16}}
  \closeproofpartwideindR

  \proofpartwide{Schluss}
    \proofstepwidestar[1]{A,B\in\mathbb R}{\rA}
    \proofstepwidestar[2]{\neg(A\subseteq B\lor B\subseteq A)}{\rA}
    \proofstepwidestar[2]{A\nsubseteq B\land B\nsubseteq A}{%
      \FormulaRefAuto{\neg(P\lor Q)\eqvdash\neg P\land\neg Q}{2}}
    \proofstepwidestar[1,2]{\exists a\in A\setminus B}{%
      \FormulaRefAuto{NonSubsetWitness}[thm]{3}}
    \proofstepwidestar[1,2]{\exists b\in B\setminus A}{%
      \FormulaRefAuto{NonSubsetWitness}[thm]{3}}
    \proofstepwidestar[1,2]{a\in A\setminus B}{\rA}
    \proofstepwidestar[1,2]{b\in B\setminus A}{\rA}
    \proofstepwidestar[1,2]{\bot}{%
      \FormulaRefAuto{RealCutOppositeWitnessesImpossible}[thm]{1,6,7}}
    \proofstepwidestar[1,2]{\bot}{\rEE{5,7,8}}
    \proofstepwidestar[1,2]{\bot}{\rEE{4,6,9}}
    \proofstepwidestar[1]{A\subseteq B\lor B\subseteq A}{\rCI{2,10}}
  \closeproofpartwide
\end{tabproofsplitwide}

\FormulaThmDeltaK[Die Inklusionsordnung der Schnitte ist total]{%
  \TotOrd{\mathbb R,\RealLe}%
}{RealCutTotalOrder}{%
  \DeltaRow{Mengen}{A\dsep B}%
  \DeltaRow{Relationsoperatoren}{\RealLe}%
}
\begin{tabproofsplitwide}
  \proofpartwideindR[Partielle Ordnung]{%
    \PartOrd{\mathbb R,\RealLe}%
  }{%
    \PartOrd{\mathbb R,\RealLe}%
  }[RealCutInclusionPartialOrder]
    \proofstepwidestarbreakkeep[1]{A\in\mathbb R}{\rA}
    \proofstepwidestarbreakkeep[1]{A\subseteq A}{%
      \FormulaRefAuto{A\subseteq A}}
    \proofstepwidestarbreakkeep[1]{A\RealLe A}{%
      \FormulaRefAuto{RealCutOrder}[def]{1,2}}
    \proofstepwidestarbreakkeep[]{%
      \forall A\in\mathbb R\,(A\RealLe A)}{\rUI{\rRI{1,3}}}
    \proofstepwidestarbreakkeep[5]{A,B,C\in\mathbb R}{\rA}
    \proofstepwidestarbreakkeep[6]{A\RealLe B\land B\RealLe C}{\rA}
    \proofstepwidestarbreakkeep[5,6]{A\subseteq B\land B\subseteq C}{%
      \FormulaRefAuto{RealCutOrderCharacterization}[thm]{5,6}}
    \proofstepwidestarbreakkeep[5,6]{A\subseteq C}{%
      \FormulaRefAuto{A\subseteq B\dsep B\subseteq C\vdash A\subseteq C}{7}}
    \proofstepwidestarbreakkeep[5,6]{A\RealLe C}{%
      \FormulaRefAuto{RealCutOrder}[def]{5,8}}
    \proofstepwidestarbreakkeep[]{%
      \forall A,B,C\in\mathbb R\,
      ((A\RealLe B\land B\RealLe C)\rightarrow A\RealLe C)}{%
      \rUI{\rRI{5,\rRI{6,9}}}}
    \proofstepwidestarbreakkeep[11]{A,B\in\mathbb R}{\rA}
    \proofstepwidestarbreakkeep[12]{A\RealLe B\land B\RealLe A}{\rA}
    \proofstepwidestarbreakkeep[11,12]{A\subseteq B\land B\subseteq A}{%
      \FormulaRefAuto{RealCutOrderCharacterization}[thm]{11,12}}
    \proofstepwidestarbreakkeep[11,12]{A=B}{%
      \FormulaRefAuto{A\subseteq B\dsep B\subseteq A\vdash A=B}{13}}
    \proofstepwidestarbreakkeep[]{%
      \forall A,B\in\mathbb R\,
      ((A\RealLe B\land B\RealLe A)\rightarrow A=B)}{%
      \rUI{\rRI{11,\rRI{12,14}}}}
    \proofstepwidestarbreakkeep[]{\PartOrd{\mathbb R,\RealLe}}{%
      \FormulaRefAuto{Partielle Ordnung}[def]{4,10,15}}
  \closeproofpartwideindR

  \proofpartwide{Vergleichbarkeit und Schluss}
    \proofstepwidestar[]{\PartOrd{\mathbb R,\RealLe}}{%
      \FormulaRefAuto{RealCutInclusionPartialOrder}[thm]}
    \proofstepwidestar[2]{A,B\in\mathbb R}{\rA}
    \proofstepwidestar[2]{A\subseteq B\lor B\subseteq A}{%
      \FormulaRefAuto{RealCutInclusionComparable}[thm]{2}}
    \proofstepwidestar[2]{A\RealLe B\lor B\RealLe A}{%
      \FormulaRefAuto{RealCutOrderCharacterization}[thm]{2,3}}
    \proofstepwidestar[]{%
      \forall A,B\in\mathbb R\,(A\RealLe B\lor B\RealLe A)}{%
      \rUI{\rRI{2,4}}}
    \proofstepwidestar[]{\TotOrd{\mathbb R,\RealLe}}{%
      \FormulaRefAuto{Totale Ordnung}[def]{1,5}}
  \closeproofpartwide
\end{tabproofsplitwide}

\FormulaThmDeltaK[Hauptschnitte bewahren und spiegeln die Ordnung]{%
  p,q\in\mathbb Q\vdash
  \bigl(p\leq q\leftrightarrow\RealEmb(p)\RealLe\RealEmb(q)\bigr)
  \land
  \bigl(p<q\leftrightarrow\RealEmb(p)\RealLt\RealEmb(q)\bigr)%
}{RationalToRealEmbeddingOrder}{%
  \DeltaRow{Mengen}{p\dsep q}%
  \DeltaRow{Funktionen}{\RealEmb}%
  \DeltaRow{Relationsoperatoren}{\RealLe\dsep\RealLt}%
}
\begin{tabproofsplitwide}
  \proofpartwideindR[Nichtstrikte Hinrichtung]{%
    p,q\in\mathbb Q\dsep p\leq q
    \vdash\RealHat p\subseteq\RealHat q
  }{%
    p,q\in\mathbb Q\dsep p\leq q
    \vdash\RealHat p\subseteq\RealHat q
  }[RealPrincipalCutOrderForward]
    \proofstepwidestar[1]{p,q\in\mathbb Q}{\rA}
    \proofstepwidestar[2]{p\leq q}{\rA}
    \proofstepwidestar[3]{x\in\RealHat p}{\rA}
    \proofstepwide[1,3]{x}{<}{p}{%
      \FormulaRefAuto{RealPrincipalCut}[def]{1,3}}
    \proofstepwide[1,2,3]{x}{<}{q}{%
      \FormulaRefAuto{RationalOrderTotal}[thm]{2,4}}
    \proofstepwidestar[1,2,3]{x\in\RealHat q}{%
      \FormulaRefAuto{RealPrincipalCut}[def]{1,5}}
    \proofstepwidestar[1,2]{\RealHat p\subseteq\RealHat q}{%
      \FormulaRefAuto{SubsetIntroductionRule}[thm]{[3]\vdots 6}}
  \closeproofpartwideindR

  \proofpartwideindR[Nichtstrikte Rückrichtung]{%
    p,q\in\mathbb Q\dsep\RealHat p\subseteq\RealHat q
    \vdash p\leq q
  }{%
    p,q\in\mathbb Q\dsep\RealHat p\subseteq\RealHat q
    \vdash p\leq q
  }[RealPrincipalCutOrderBackward]
    \proofstepwidestar[1]{p,q\in\mathbb Q}{\rA}
    \proofstepwidestar[2]{\RealHat p\subseteq\RealHat q}{\rA}
    \proofstepwidestar[3]{q<p}{\rA}
    \proofstepwidestar[1,3]{%
      \exists r\in\mathbb Q\,(q<r\land r<p)}{%
      \FormulaRefAuto{RationalOrderDense}[thm]{1,3}}
    \proofstepwidestar[1,3]{r\in\mathbb Q\land q<r\land r<p}{\rA}
    \proofstepwidestar[1,3]{r\in\RealHat p}{%
      \FormulaRefAuto{RealPrincipalCut}[def]{5}}
    \proofstepwidestar[1,2,3]{r\in\RealHat q}{%
      \FormulaRefAuto{A\subseteq B\dsep x\in A\vdash x\in B}{2,6}}
    \proofstepwidestar[1,3]{r\notin\RealHat q}{%
      \FormulaRefAuto{RealPrincipalCut}[def]{5}}
    \proofstepwidestar[1,2,3]{\bot}{\rBI{7,8}}
    \proofstepwidestar[1,2,3]{\bot}{\rEE{4,5,9}}
    \proofstepwidestar[1,2]{\neg(q<p)}{\rCI{3,10}}
    \proofstepwidestar[1,2]{p\leq q}{%
      \FormulaRefAuto{RationalOrderTotal}[thm]{1,11}}
  \closeproofpartwideindR

  \proofpartwideindR[Strikte Hinrichtung]{%
    p,q\in\mathbb Q\dsep p<q
    \vdash\RealHat p\subsetneq\RealHat q
  }{%
    p,q\in\mathbb Q\dsep p<q
    \vdash\RealHat p\subsetneq\RealHat q
  }[RealPrincipalCutStrictOrderForward]
    \proofstepwidestar[1]{p,q\in\mathbb Q}{\rA}
    \proofstepwidestar[2]{p<q}{\rA}
    \proofstepwidestar[1,2]{p\leq q}{%
      \FormulaRefAuto{RationalStrictOrderDef}[def]{1,2}}
    \proofstepwidestar[1,2]{\RealHat p\subseteq\RealHat q}{%
      \FormulaRefAuto{RealPrincipalCutOrderForward}[thm]{1,3}}
    \proofstepwidestar[1,2]{%
      \exists r\in\mathbb Q\,(p<r\land r<q)}{%
      \FormulaRefAuto{RationalOrderDense}[thm]{1,2}}
    \proofstepwidestar[6]{r\in\mathbb Q\land p<r\land r<q}{\rA}
    \proofstepwidestar[1,2,6]{r\in\RealHat q}{%
      \FormulaRefAuto{RealPrincipalCut}[def]{6}}
    \proofstepwidestar[1,2,6]{r\notin\RealHat p}{%
      \FormulaRefAuto{RealPrincipalCut}[def]{6}}
    \proofstepwidestar[1,2,6]{\RealHat p\neq\RealHat q}{%
      \FormulaRefAuto{x\in B, x\not\in A\vdash A\neq B}{7,8}}
    \proofstepwidestar[1,2,6]{\RealHat p\subsetneq\RealHat q}{%
      \FormulaRefAuto{A\subset B\coloneq A\subseteq B\land A\neq B}{4,9}}
    \proofstepwidestar[1,2]{\RealHat p\subsetneq\RealHat q}{%
      \rEE{5,6,10}}
  \closeproofpartwideindR

  \proofpartwideindR[Strikte Rückrichtung]{%
    p,q\in\mathbb Q\dsep\RealHat p\subsetneq\RealHat q
    \vdash p<q
  }{%
    p,q\in\mathbb Q\dsep\RealHat p\subsetneq\RealHat q
    \vdash p<q
  }[RealPrincipalCutStrictOrderBackward]
    \proofstepwidestar[1]{p,q\in\mathbb Q}{\rA}
    \proofstepwidestar[2]{\RealHat p\subsetneq\RealHat q}{\rA}
    \proofstepwidestar[2]{%
      \RealHat p\subseteq\RealHat q\land\RealHat p\neq\RealHat q}{%
      \FormulaRefAuto{A\subset B\coloneq A\subseteq B\land A\neq B}{2}}
    \proofstepwidestar[1,2]{p\leq q}{%
      \FormulaRefAuto{RealPrincipalCutOrderBackward}[thm]{1,3}}
    \proofstepwidestar[5]{p=q}{\rA}
    \proofstepwidestar[1,5]{\RealHat p=\RealHat q}{\rIE{5}}
    \proofstepwidestar[1,2,5]{\bot}{\rBI{3,6}}
    \proofstepwidestar[1,2]{p\neq q}{\rCI{5,7}}
    \proofstepwidestar[1,2]{p<q}{%
      \FormulaRefAuto{RationalStrictOrderDef}[def]{1,4,8}}
  \closeproofpartwideindR

  \proofpartwide{Einbettung und Zusammenfassung}
    \proofstepwidestarbreakkeep[1]{p,q\in\mathbb Q}{\rA}
    \proofstepwidestar[1]{p,q\in\mathbb Q\dsep p\leq q\vdash\RealHat p\subseteq\RealHat q}{%
      \FormulaRefAuto{RealPrincipalCutOrderForward}[thm]}
    \proofstepwidestar[1]{p,q\in\mathbb Q\dsep\RealHat p\subseteq\RealHat q\vdash p\leq q}{%
      \FormulaRefAuto{RealPrincipalCutOrderBackward}[thm]}
    \proofstepwide[1]{p\leq q}{\leftrightarrow}{%
      \RealHat p\subseteq\RealHat q}{%
      \rLRI{2,
        3}}
    \proofstepwidestar[1]{p,q\in\mathbb Q\dsep p<q\vdash\RealHat p\subsetneq\RealHat q}{%
      \FormulaRefAuto{RealPrincipalCutStrictOrderForward}[thm]}
    \proofstepwidestar[1]{p,q\in\mathbb Q\dsep\RealHat p\subsetneq\RealHat q\vdash p<q}{%
      \FormulaRefAuto{RealPrincipalCutStrictOrderBackward}[thm]}
    \proofstepwide[1]{p<q}{\leftrightarrow}{%
      \RealHat p\subsetneq\RealHat q}{%
      \rLRI{5,
        6}}
    \proofstepwidestarbreakkeep[1]{%
      \RealEmb(p)=\RealHat p\land\RealEmb(q)=\RealHat q}{%
      \FormulaRefAuto{RationalToRealEmbedding}[def]{1}}
    \proofstepwide[1]{p\leq q}{\leftrightarrow}{%
      \RealEmb(p)\RealLe\RealEmb(q)}{%
      \FormulaRefAuto{RealCutOrderCharacterization}[thm]{4,8}}
    \proofstepwide[1]{p<q}{\leftrightarrow}{%
      \RealEmb(p)\RealLt\RealEmb(q)}{%
      \FormulaRefAuto{RealCutOrderCharacterization}[thm]{7,8}}
    \proofstepwidestarbreakkeep[1]{%
      \begin{aligned}[t]
        &(p\leq q\leftrightarrow\RealEmb(p)\RealLe\RealEmb(q))\land{}\\[-2pt]
        &(p<q\leftrightarrow\RealEmb(p)\RealLt\RealEmb(q))
      \end{aligned}}{%
      \rAI{9,10}}
  \closeproofpartwide
\end{tabproofsplitwide}

\FormulaThmDeltaK[Injektivität der rationalen Einbettung]{%
  \RealEmb\colon\mathbb Q\inj\mathbb R%
}{RationalToRealEmbeddingInjective}{%
  \DeltaRow{Funktionen}{\RealEmb}%
}
\begin{tabproofsplitwide}
  \proofpartwideindR[Punktweise Injektivität]{%
    p,q\in\mathbb Q\dsep\RealEmb(p)=\RealEmb(q)\vdash p=q
  }{%
    p,q\in\mathbb Q\dsep\RealEmb(p)=\RealEmb(q)\vdash p=q
  }[RationalToRealEmbeddingPointInjective]
    \proofstepwidestar[1]{p,q\in\mathbb Q}{\rA}
    \proofstepwidestar[2]{\RealEmb(p)=\RealEmb(q)}{\rA}
    \proofstepwidestar[1,2]{%
      \RealEmb(p)\RealLe\RealEmb(q)\land
      \RealEmb(q)\RealLe\RealEmb(p)}{%
      \FormulaRefAuto{RealCutOrderCharacterization}[thm]{1,2}}
    \proofstepwidestar[1,2]{p\leq q\land q\leq p}{%
      \FormulaRefAuto{RationalToRealEmbeddingOrder}[thm]{1,3}}
    \proofstepwidestar[1,2]{p=q}{%
      \FormulaRefAuto{RationalOrderTotal}[thm]{1,4}}
  \closeproofpartwideindR

  \proofpartwide{Funktionsschluss}
    \proofstepwidestar[]{\RealEmb\colon\mathbb Q\to\mathbb R}{%
      \FormulaRefAuto{RationalToRealEmbeddingFunction}[thm]}
    \proofstepwidestar[]{%
      \forall p,q\in\mathbb Q\,
      (\RealEmb(p)=\RealEmb(q)\rightarrow p=q)}{%
      \FormulaRefAuto{RationalToRealEmbeddingPointInjective}[thm]}
    \proofstepwidestar[]{\RealEmb\colon\mathbb Q\inj\mathbb R}{%
      \FormulaRefAuto{Injektive Funktion}[def]{1,2}}
  \closeproofpartwide
\end{tabproofsplitwide}


% Abschnitt 3.2 des Hauptbands (Zahlbereichskopien) gehört nicht zu diesem Kern.
\setcounter{section}{2}
% Kanonische Quelle: Band 19 und Beweistabellen verwenden denselben Text.
\section{Suprema als Vereinigungen}

\FormulaThmDeltaK[Vereinigung einer beschränkten Schnittfamilie]{%
  \begin{aligned}[t]
    &\varnothing\neq\mathcal A\subseteq\mathbb R\dsep
      U\in\mathbb R\dsep
      \forall A\in\mathcal A\,(A\RealLe U)\\[-2pt]
    &\hspace{34mm}\vdash\RealCut{\bigcup\mathcal A}.
  \end{aligned}%
}{RealBoundedCutUnionIsCut}{%
  \DeltaRow{Mengen}{\mathcal A\dsep A\dsep U}%
  \DeltaRow{Relationsoperatoren}{\RealLe}%
}
\begin{tabproofsplitwide}
\proofpartwide{Nichtleerheit}
\proofstepwidestarbreakkeep[1]{\varnothing\neq\mathcal A\subseteq\mathbb R}{\rA}
    \proofstepwidestarbreakkeep[2]{U\in\mathbb R}{\rA}
    \proofstepwidestarbreakkeep[3]{\forall A\in\mathcal A\,(A\RealLe U)}{\rA}
    \proofstepwidestarbreakkeep[1]{\exists A_0\in\mathcal A}{%
      \FormulaRefAuto{A\neq\varnothing\eqvdash\exists x\,(x\in A)}{1}}
    \proofstepwidestarbreakkeep[1]{A_0\in\mathcal A}{\rA}
    \proofstepwidestarbreakkeep[1,5]{A_0\in\mathbb R}{%
      \FormulaRefAuto{A\subseteq B\dsep x\in A\vdash x\in B}{1,5}}
    \proofstepwidestarbreakkeep[1,5]{\RealCut{A_0}}{%
      \FormulaRefAuto{RealCutMembership}[thm]{6}}
    \proofstepwidestarbreakkeep[1,5]{\exists a_0\in A_0}{%
      \FormulaRefAuto{DedekindCutCharacterization}[thm]{7}}
    \proofstepwidestarbreakkeep[1,5]{a_0\in A_0}{\rA}
    \proofstepwidestarbreakkeep[1,5]{a_0\in\bigcup\mathcal A}{%
      \FormulaRefAuto{FamilyMemberSubsetUnion}[thm]{5,9}}
    \proofstepwidestarbreakkeep[1,5]{\bigcup\mathcal A\neq\varnothing}{%
      \FormulaRefAuto{x\in A\vdash A\neq\varnothing}{10}}
    \proofstepwidestarbreakkeep[1]{\bigcup\mathcal A\neq\varnothing}{%
      \rEE{8,9,11}}
    \proofstepwidestarbreakkeep[1]{\bigcup\mathcal A\neq\varnothing}{%
      \rEE{4,5,12}}

\closeproofpartwide

\proofpartwide{Trägerinklusion}
  \setcounter{proofstepnr}{13}% Fortlaufende Schritte dieses Beweises.
\proofstepwidestarbreakkeep[14]{x\in\bigcup\mathcal A}{\rA}
    \proofstepwidestarbreakkeep[14]{\exists A\in\mathcal A\,(x\in A)}{%
      \FormulaRefAuto{x\in\bigcup A\eqvdash\exists B\in A\,(x\in B)}{14}}
    \proofstepwidestarbreakkeep[14]{A\in\mathcal A\land x\in A}{\rA}
    \proofstepwidestarbreakkeep[1,14]{A\in\mathbb R}{%
      \FormulaRefAuto{A\subseteq B\dsep x\in A\vdash x\in B}{1,16}}
    \proofstepwidestarbreakkeep[1,14]{A\subseteq\mathbb Q}{%
      \FormulaRefAuto{RealCutMembership}[thm]{17}}
    \proofstepwidestarbreakkeep[1,14]{x\in\mathbb Q}{%
      \FormulaRefAuto{A\subseteq B\dsep x\in A\vdash x\in B}{18,16}}
    \proofstepwidestarbreakkeep[1,14]{x\in\mathbb Q}{\rEE{15,16,19}}
    \proofstepwidestarbreakkeep[1]{\bigcup\mathcal A\subseteq\mathbb Q}{%
      \FormulaRefAuto{SubsetIntroductionRule}[thm]{[14]\vdots 20}}

\closeproofpartwide

\proofpartwide{Echtheit}
  \setcounter{proofstepnr}{21}% Fortlaufende Schritte dieses Beweises.
\proofstepwidestarbreakkeep[22]{x\in\bigcup\mathcal A}{\rA}
    \proofstepwidestarbreakkeep[22]{\exists A\in\mathcal A\,(x\in A)}{%
      \FormulaRefAuto{x\in\bigcup A\eqvdash\exists B\in A\,(x\in B)}{22}}
    \proofstepwidestarbreakkeep[22]{A\in\mathcal A\land x\in A}{\rA}
    \proofstepwidestarbreakkeep[3,22]{A\RealLe U}{\rUE{3,24}}
    \proofstepwidestarbreakkeep[3,22]{A\subseteq U}{%
      \FormulaRefAuto{RealCutOrderCharacterization}[thm]{2,24,25}}
    \proofstepwidestarbreakkeep[3,22]{x\in U}{%
      \FormulaRefAuto{A\subseteq B\dsep x\in A\vdash x\in B}{26,24}}
    \proofstepwidestarbreakkeep[3,22]{x\in U}{\rEE{23,24,27}}
    \proofstepwidestarbreakkeep[3]{\bigcup\mathcal A\subseteq U}{%
      \FormulaRefAuto{SubsetIntroductionRule}[thm]{[22]\vdots 28}}
    \proofstepwidestarbreakkeep[2]{\RealCut U}{%
      \FormulaRefAuto{RealCutMembership}[thm]{2}}
    \proofstepwidestarbreakkeep[2]{U\neq\mathbb Q}{%
      \FormulaRefAuto{DedekindCutCharacterization}[thm]{30}}
    \proofstepwidestarbreakkeep[32]{\bigcup\mathcal A=\mathbb Q}{\rA}
    \proofstepwidestarbreakkeep[2,3,32]{\mathbb Q\subseteq U}{\rIE{32,29}}
    \proofstepwidestarbreakkeep[2]{U\subseteq\mathbb Q}{%
      \FormulaRefAuto{DedekindCutCharacterization}[thm]{30}}
    \proofstepwidestarbreakkeep[2,3,32]{U=\mathbb Q}{%
      \FormulaRefAuto{A\subseteq B\dsep B\subseteq A\vdash A=B}{34,33}}
    \proofstepwidestarbreakkeep[2,3,32]{\bot}{\rBI{31,35}}
    \proofstepwidestarbreakkeep[2,3]{\bigcup\mathcal A\neq\mathbb Q}{%
      \rCI{32,36}}

\closeproofpartwide

\proofpartwideindR[Träger, Nichtleerheit und Echtheit]{%
    \begin{aligned}[t]
      &\varnothing\neq\mathcal A\subseteq\mathbb R\dsep U\in\mathbb R\dsep
        \forall A\in\mathcal A\,(A\RealLe U)\\[-2pt]
      &\quad\vdash
        \bigcup\mathcal A\subseteq\mathbb Q\land
        \bigcup\mathcal A\neq\varnothing\land
        \bigcup\mathcal A\neq\mathbb Q
    \end{aligned}
  }{%
    \varnothing\neq\mathcal A\subseteq\mathbb R\dsep U\in\mathbb R\dsep
    \forall A\in\mathcal A\,(A\RealLe U)
    \vdash
    \bigcup\mathcal A\subseteq\mathbb Q\land
    \bigcup\mathcal A\neq\varnothing\land
    \bigcup\mathcal A\neq\mathbb Q
  }[RealBoundedCutUnionCarrierProper]
  \setcounter{proofstepnr}{37}% Fortlaufende Schritte dieses Beweises.
\proofstepwidestarbreakkeep[1,2,3]{%
      \bigcup\mathcal A\subseteq\mathbb Q\land
      \bigcup\mathcal A\neq\varnothing\land
      \bigcup\mathcal A\neq\mathbb Q}{%
      \rAI{21,\rAI{13,37}}}

\closeproofpartwideindR


  \proofpartwideindR[Abwärtsabschluss]{%
    \begin{aligned}[t]
      &\mathcal A\subseteq\mathbb R\dsep p\in\bigcup\mathcal A\dsep
        q\in\mathbb Q\dsep q<p\\[-2pt]
      &\quad\vdash q\in\bigcup\mathcal A
    \end{aligned}
  }{%
    \mathcal A\subseteq\mathbb R\dsep p\in\bigcup\mathcal A\dsep
    q\in\mathbb Q\dsep q<p\vdash q\in\bigcup\mathcal A
  }[RealCutUnionDownwardClosed]
    \proofstepwidestar[1]{\mathcal A\subseteq\mathbb R}{\rA}
    \proofstepwidestar[2]{p\in\bigcup\mathcal A}{\rA}
    \proofstepwidestar[3]{q\in\mathbb Q}{\rA}
    \proofstepwidestar[4]{q<p}{\rA}
    \proofstepwidestar[2]{\exists A\in\mathcal A\,(p\in A)}{%
      \FormulaRefAuto{x\in\bigcup A\eqvdash\exists B\in A\,(x\in B)}{2}}
    \proofstepwidestar[2]{A\in\mathcal A\land p\in A}{\rA}
    \proofstepwidestar[1,2]{A\in\mathbb R}{%
      \FormulaRefAuto{A\subseteq B\dsep x\in A\vdash x\in B}{1,6}}
    \proofstepwidestar[1,2]{\RealCut A}{%
      \FormulaRefAuto{RealCutMembership}[thm]{7}}
    \proofstepwidestar[1,2,3,4]{q\in A}{%
      \FormulaRefAuto{DedekindCutCharacterization}[thm]{8,6,3,4}}
    \proofstepwidestar[1,2,3,4]{q\in\bigcup\mathcal A}{%
      \FormulaRefAuto{FamilyMemberSubsetUnion}[thm]{6,9}}
    \proofstepwidestar[1,2,3,4]{q\in\bigcup\mathcal A}{%
      \rEE{5,6,10}}
  \closeproofpartwideindR

  \proofpartwideindR[Kein größtes Element]{%
    \mathcal A\subseteq\mathbb R\dsep p\in\bigcup\mathcal A
    \vdash\exists r\in\bigcup\mathcal A\,(p<r)
  }{%
    \mathcal A\subseteq\mathbb R\dsep p\in\bigcup\mathcal A
    \vdash\exists r\in\bigcup\mathcal A\,(p<r)
  }[RealCutUnionNoMaximum]
    \proofstepwidestar[1]{\mathcal A\subseteq\mathbb R}{\rA}
    \proofstepwidestar[2]{p\in\bigcup\mathcal A}{\rA}
    \proofstepwidestar[2]{\exists A\in\mathcal A\,(p\in A)}{%
      \FormulaRefAuto{x\in\bigcup A\eqvdash\exists B\in A\,(x\in B)}{2}}
    \proofstepwidestar[2]{A\in\mathcal A\land p\in A}{\rA}
    \proofstepwidestar[1,2]{A\in\mathbb R}{%
      \FormulaRefAuto{A\subseteq B\dsep x\in A\vdash x\in B}{1,4}}
    \proofstepwidestar[1,2]{\RealCut A}{%
      \FormulaRefAuto{RealCutMembership}[thm]{5}}
    \proofstepwidestar[1,2]{\exists r\in A\,(p<r)}{%
      \FormulaRefAuto{DedekindCutCharacterization}[thm]{6,4}}
    \proofstepwidestar[1,2]{r\in A\land p<r}{\rA}
    \proofstepwidestar[1,2]{r\in\bigcup\mathcal A}{%
      \FormulaRefAuto{FamilyMemberSubsetUnion}[thm]{4,8}}
    \proofstepwidestar[1,2]{\exists r\in\bigcup\mathcal A\,(p<r)}{%
      \rEI{8,9}}
    \proofstepwidestar[1,2]{\exists r\in\bigcup\mathcal A\,(p<r)}{%
      \rEE{7,8,10}}
    \proofstepwidestar[1,2]{\exists r\in\bigcup\mathcal A\,(p<r)}{%
      \rEE{3,4,11}}
  \closeproofpartwideindR

  \proofpartwide{Schluss}
    \proofstepwidestar[1]{\varnothing\neq\mathcal A\subseteq\mathbb R}{\rA}
    \proofstepwidestar[2]{U\in\mathbb R}{\rA}
    \proofstepwidestar[3]{\forall A\in\mathcal A\,(A\RealLe U)}{\rA}
    \proofstepwidestar[1,2,3]{%
      \bigcup\mathcal A\subseteq\mathbb Q\land
      \bigcup\mathcal A\neq\varnothing\land
      \bigcup\mathcal A\neq\mathbb Q}{%
      \FormulaRefAuto{RealBoundedCutUnionCarrierProper}[thm]{1,2,3}}
    \proofstepwidestar[1]{%
      \forall p\in\bigcup\mathcal A\,\forall q\in\mathbb Q\,
      (q<p\rightarrow q\in\bigcup\mathcal A)}{%
      \FormulaRefAuto{RealCutUnionDownwardClosed}[thm]{1}}
    \proofstepwidestar[1]{%
      \forall p\in\bigcup\mathcal A\,
      \exists r\in\bigcup\mathcal A\,(p<r)}{%
      \FormulaRefAuto{RealCutUnionNoMaximum}[thm]{1}}
    \proofstepwidestar[1,2,3]{\RealCut{\bigcup\mathcal A}}{%
      \FormulaRefAuto{DedekindCutCharacterization}[thm]{4,5,6}}
  \closeproofpartwide
\end{tabproofsplitwide}

\FormulaThmDeltaK[Vollständigkeit der reellen Ordnung]{%
  \begin{aligned}[t]
    &\varnothing\neq\mathcal A\subseteq\mathbb R\dsep
      \exists U\in\mathbb R\,
        \forall A\in\mathcal A\,(A\RealLe U)\\[-2pt]
    &\hspace{19mm}\vdash
      \exists S\in\mathbb R\,
      \left(
        \begin{aligned}[t]
        &\forall A\in\mathcal A\,(A\RealLe S)\land{}\\[-2pt]
        &\forall V\in\mathbb R\,
          \bigl(\forall A\in\mathcal A\,(A\RealLe V)
          \rightarrow S\RealLe V\bigr)
        \end{aligned}
      \right).
  \end{aligned}%
}{RealLeastUpperBoundProperty}{%
  \DeltaRow{Mengen}{\mathcal A\dsep A\dsep S\dsep U\dsep V}%
  \DeltaRow{Relationsoperatoren}{\RealLe}%
}
\begin{tabproofsplitwide}
  \proofpartwideindR[Vereinigung ist obere Schranke]{%
    \begin{aligned}[t]
      &\varnothing\neq\mathcal A\subseteq\mathbb R\dsep
        S=\bigcup\mathcal A\dsep S\in\mathbb R\\[-2pt]
      &\quad\vdash\forall A\in\mathcal A\,(A\RealLe S)
    \end{aligned}
  }{%
    \varnothing\neq\mathcal A\subseteq\mathbb R\dsep
    S=\bigcup\mathcal A\dsep S\in\mathbb R
    \vdash\forall A\in\mathcal A\,(A\RealLe S)
  }[RealCutUnionUpperBound]
    \proofstepwidestar[1]{\varnothing\neq\mathcal A\subseteq\mathbb R}{\rA}
    \proofstepwidestar[2]{S=\bigcup\mathcal A}{\rA}
    \proofstepwidestar[3]{S\in\mathbb R}{\rA}
    \proofstepwidestar[4]{A\in\mathcal A}{\rA}
    \proofstepwidestar[1,4]{A\in\mathbb R}{%
      \FormulaRefAuto{A\subseteq B\dsep x\in A\vdash x\in B}{1,4}}
    \proofstepwidestar[4]{A\subseteq\bigcup\mathcal A}{%
      \FormulaRefAuto{FamilyMemberSubsetUnion}[thm]{4}}
    \proofstepwidestar[2,4]{A\subseteq S}{\rIE{2,6}}
    \proofstepwidestar[1,2,3,4]{A\RealLe S}{%
      \FormulaRefAuto{RealCutOrder}[def]{5,3,7}}
    \proofstepwidestar[1,2,3]{\forall A\in\mathcal A\,(A\RealLe S)}{%
      \rUI{\rRI{4,8}}}
  \closeproofpartwideindR

  \proofpartwideindR[Vereinigung liegt unter jeder oberen Schranke]{%
    \begin{aligned}[t]
      &\mathcal A\subseteq\mathbb R\dsep\bigcup\mathcal A\in\mathbb R\\[-2pt]
      &\quad\vdash\forall V\in\mathbb R\,
        \bigl(\forall A\in\mathcal A\,(A\RealLe V)
        \rightarrow\bigcup\mathcal A\RealLe V\bigr)
    \end{aligned}
  }{%
    \mathcal A\subseteq\mathbb R\dsep\bigcup\mathcal A\in\mathbb R
    \vdash\forall V\in\mathbb R\,
    (\forall A\in\mathcal A\,(A\RealLe V)
    \rightarrow\bigcup\mathcal A\RealLe V)
  }[RealCutUnionLeastUpperBound]
    \proofstepwidestar[1]{\mathcal A\subseteq\mathbb R}{\rA}
    \proofstepwidestar[2]{\bigcup\mathcal A\in\mathbb R}{\rA}
    \proofstepwidestar[3]{V\in\mathbb R}{\rA}
    \proofstepwidestar[4]{\forall A\in\mathcal A\,(A\RealLe V)}{\rA}
    \proofstepwidestar[5]{x\in\bigcup\mathcal A}{\rA}
    \proofstepwidestar[5]{\exists A\in\mathcal A\,(x\in A)}{%
      \FormulaRefAuto{x\in\bigcup A\eqvdash\exists B\in A\,(x\in B)}{5}}
    \proofstepwidestar[7]{A\in\mathcal A\land x\in A}{\rA}
    \proofstepwidestar[1,7]{A\in\mathbb R}{%
      \FormulaRefAuto{A\subseteq B\dsep x\in A\vdash x\in B}{1,7}}
    \proofstepwidestar[4,7]{A\RealLe V}{\rUE{4,7}}
    \proofstepwidestar[1,3,4,7]{A\subseteq V}{%
      \FormulaRefAuto{RealCutOrderCharacterization}[thm]{8,3,9}}
    \proofstepwidestar[1,3,4,7]{x\in V}{%
      \FormulaRefAuto{A\subseteq B\dsep x\in A\vdash x\in B}{10,7}}
    \proofstepwidestar[1,3,4,5]{x\in V}{\rEE{6,7,11}}
    \proofstepwidestar[1,3,4]{\bigcup\mathcal A\subseteq V}{%
      \FormulaRefAuto{SubsetIntroductionRule}[thm]{[5]\vdots 12}}
    \proofstepwidestar[1,2,3,4]{\bigcup\mathcal A\RealLe V}{%
      \FormulaRefAuto{RealCutOrder}[def]{2,3,13}}
    \proofstepwidestar[1,2,3]{%
      \forall A\in\mathcal A\,(A\RealLe V)
      \rightarrow\bigcup\mathcal A\RealLe V}{\rRI{4,14}}
    \proofstepwidestar[1,2]{%
      \forall V\in\mathbb R\,
      (\forall A\in\mathcal A\,(A\RealLe V)
       \rightarrow\bigcup\mathcal A\RealLe V)}{\rUI{\rRI{3,15}}}
  \closeproofpartwideindR

  \proofpartwide{Existenz der kleinsten oberen Schranke}
    \proofstepwidestarbreakkeep[1]{\varnothing\neq\mathcal A\subseteq\mathbb R}{\rA}
    \proofstepwidestarbreakkeep[2]{%
      \exists U\in\mathbb R\,\forall A\in\mathcal A\,(A\RealLe U)}{\rA}
    \proofstepwidestarbreakkeep[2]{%
      U\in\mathbb R\land\forall A\in\mathcal A\,(A\RealLe U)}{\rA}
    \proofstepwidestarbreakkeep[1,2]{\RealCut{\bigcup\mathcal A}}{%
      \FormulaRefAuto{RealBoundedCutUnionIsCut}[thm]{1,3}}
    \proofstepwidestarbreakkeep[1,2]{\bigcup\mathcal A\in\mathbb R}{%
      \FormulaRefAuto{RealCutMembership}[thm]{4}}
    \proofstepwidestarbreakkeep[1,2]{%
      \forall A\in\mathcal A\,(A\RealLe\bigcup\mathcal A)}{%
      \FormulaRefAuto{RealCutUnionUpperBound}[thm]{1,\rII,5}}
    \proofstepwidestarbreakkeep[1,2]{%
      \forall V\in\mathbb R\,
      (\forall A\in\mathcal A\,(A\RealLe V)
       \rightarrow\bigcup\mathcal A\RealLe V)}{%
      \FormulaRefAuto{RealCutUnionLeastUpperBound}[thm]{1,5}}
    \proofstepwidestarbreakkeep[1,2]{%
      \exists S\in\mathbb R\,
      (\forall A\in\mathcal A\,(A\RealLe S)\land
       \forall V\in\mathbb R\,
       (\forall A\in\mathcal A\,(A\RealLe V)\rightarrow S\RealLe V))}{%
      \rEI{5,\rAI{6,7}}}
    \proofstepwidestarbreakkeep[1,2]{%
      \exists S\in\mathbb R\,
      (\forall A\in\mathcal A\,(A\RealLe S)\land
       \forall V\in\mathbb R\,
       (\forall A\in\mathcal A\,(A\RealLe V)\rightarrow S\RealLe V))}{%
      \rEE{2,3,8}}
  \closeproofpartwide
\end{tabproofsplitwide}

Erst die vorstehende Existenz- und Kleinste-Obergrenzen-Aussage erfüllt die
Definitionsprämisse des allgemeinen Supremumsbegriffs aus Band~13. Daher ist
der folgende Supremumsterm wohldefiniert.

\FormulaThmDeltaK[Supremum einer beschränkten Schnittfamilie]{%
  \begin{aligned}[t]
    &\varnothing\neq\mathcal A\subseteq\mathbb R\dsep
      U\in\mathbb R\dsep
      \forall A\in\mathcal A\,(A\RealLe U)\\[-2pt]
    &\hspace{20mm}\vdash
      \sup_{\mathbb R,\RealLe}(\mathcal A)=\bigcup\mathcal A.
  \end{aligned}%
}{RealCutSupremumUnion}{%
  \DeltaRow{Mengen}{\mathcal A\dsep A\dsep U}%
  \DeltaRow{Relationsoperatoren}{\RealLe}%
}
\begin{tabproofwide}
  \proofstepwidestarbreakkeep[1]{\varnothing\neq\mathcal A\subseteq\mathbb R}{\rA}
  \proofstepwidestarbreakkeep[2]{U\in\mathbb R}{\rA}
  \proofstepwidestarbreakkeep[3]{\forall A\in\mathcal A\,(A\RealLe U)}{\rA}
  \proofstepwidestar[1,2,3]{\RealCut { \bigcup \mathcal A }}{%
      \FormulaRefAuto{RealBoundedCutUnionIsCut}[thm]{1,2,3}}
  \proofstepwidestarbreakkeep[1,2,3]{\bigcup\mathcal A\in\mathbb R}{%
    \FormulaRefAuto{RealCutMembership}[thm]{%
      4}}
  \proofstepwidestarbreakkeep[1,2,3]{%
    \forall A\in\mathcal A\,(A\RealLe\bigcup\mathcal A)}{%
    \FormulaRefAuto{RealCutUnionUpperBound}[thm]{1,\rII,5}}
  \proofstepwidestarbreakkeep[1,2,3]{%
    \forall V\in\mathbb R\,
    (\forall A\in\mathcal A\,(A\RealLe V)
      \rightarrow\bigcup\mathcal A\RealLe V)}{%
    \FormulaRefAuto{RealCutUnionLeastUpperBound}[thm]{1,5}}
  \proofstepwidestarbreakkeep[1,2,3]{%
    \bigcup\mathcal A\in\UB_{\mathbb R,\RealLe}(\mathcal A)}{%
    \FormulaRefAuto{UpperBoundSetIntro}[thm]{5,6}}
  \proofstepwidestarbreakkeep[1,2,3]{%
    \forall V\in\UB_{\mathbb R,\RealLe}(\mathcal A)\,
      (\bigcup\mathcal A\RealLe V)}{%
    \FormulaRefAuto{UpperBoundSetElim}[thm]{7}}
  \proofstepwidestarbreakkeep[1,2,3]{%
    \bigcup\mathcal A=\sup_{\mathbb R,\RealLe}(\mathcal A)}{%
    \FormulaRefAuto{SupremumCandidateEqualsSupremum}[thm]{1,8,9}}
  \proofstepwidestarbreakkeep[1,2,3]{%
    \sup_{\mathbb R,\RealLe}(\mathcal A)=\bigcup\mathcal A}{%
    \FormulaRefAuto{a=b\vdash b=a}{10}}
\end{tabproofwide}


